% 山东科技大学硕士研究生学位论文开题报告 —— 封面 / 评分表 / 评审意见表
% 复刻学校 docx 制式模板版式（XeLaTeX 编译，正文见 main.tex，合并后封面在前）
\documentclass[a4paper,12pt]{ctexart}

\usepackage[margin=2.5cm]{geometry}
\usepackage{multirow, array}
\usepackage{graphicx}

\pagestyle{empty}

\begin{document}

% ==================== 封面 ====================
\thispagestyle{empty}
\begin{center}
  \vspace*{0.4cm}
  {\zihao{1}\heiti 山东科技大学}\\[10pt]
  {\zihao{2}\heiti 硕士研究生学位论文开题报告}\\[6pt]
  {\zihao{2}\heiti 及论文工作计划}
\end{center}

\vspace{1.0cm}

\begin{tabular}{@{}p{3.6cm}p{10.4cm}@{}}
  \textbf{论文题目} & 面向发动机燃烧场多组分浓度与关键自由基二维分布同步诊断的激光光谱方法研究 \\
  \cline{2-2}
  \textbf{研究生姓名} & 张洪民 \\
  \cline{2-2}
  \textbf{导师姓名、职称} & 刘训臣~教授 \\
  \cline{2-2}
  \textbf{学科、专业} & 机械 \\
  \cline{2-2}
  \textbf{研究方向} & 燃烧激光光谱诊断 \\
  \cline{2-2}
  \textbf{入学年月} & 2024年9月 \\
  \cline{2-2}
  \textbf{所 属 学 院} & 机械电子工程学院 \\
  \cline{2-2}
  \textbf{开题时间} & 2026年1月24日 \\
  \cline{2-2}
\end{tabular}

\vspace{1.2cm}

\begin{center}
  {\zihao{4}\heiti 山东科技大学学位评定委员会办公室制表}
\end{center}

\newpage

% ==================== 评分表 ====================
\begin{center}
  {\zihao{4}\heiti 山东科技大学硕士研究生学位论文开题报告评分表}
  \vspace{8pt}

  \zihao{5}
  \begin{tabular}{|p{2.6cm}|p{1.4cm}|p{2.0cm}|p{6.2cm}|p{1.6cm}|}
    \hline
    评审项目 & 分值 & \multicolumn{2}{|c|}{评 分 标 准} & 得分（百分制） \\
    \hline
    \multirow{3}{*}{一、选题依据} & \multirow{3}{*}{30分} & 21$\sim$30分 & 选题有较强的理论意义、实用价值，深刻的学术研究内涵。 & \\
    \cline{3-5}
    & & 11$\sim$20分 & 选题有一定的理论意义、实用价值，有一定的学术研究内涵。 & \\
    \cline{3-5}
    & & 0$\sim$10分 & 选题缺乏理论意义和实用价值。 & \\
    \hline
    \multirow{3}{*}{二、理论基础和专门知识} & \multirow{3}{*}{20分} & 15$\sim$20分 & 较好的掌握坚实宽广的理论基础和系统专业知识 & \\
    \cline{3-5}
    & & 6$\sim$15分 & 基本的掌握坚实宽广的理论基础和系统专业知识 & \\
    \cline{3-5}
    & & 5分以下 & 未能掌握坚实宽广的理论基础和系统知识 & \\
    \hline
    \multirow{3}{*}{三、选题难度及先进性} & \multirow{3}{*}{30分} & 21$\sim$30分 & 研究课题属本学科发展方向并居前沿位置，具有自己独特的思考、研究课题具有较强的先进性 & \\
    \cline{3-5}
    & & 11$\sim$20分 & 研究课题属本学科的发展方向，并具有先进性。 & \\
    \cline{3-5}
    & & 0$\sim$10分 & 研究课题与本学科的发展方向先进性不明显，难度欠佳。 & \\
    \hline
    \multirow{3}{*}{四、文字表达} & \multirow{3}{*}{10分} & 8$\sim$10分 & 条理清晰，分析严谨，文笔流畅 & \\
    \cline{3-5}
    & & 4$\sim$7分 & 条理较好，层次分明，文笔较流畅 & \\
    \cline{3-5}
    & & 3分以下 & 写作能力较差 & \\
    \hline
    \multirow{3}{*}{五、口头报告} & \multirow{3}{*}{10分} & 8$\sim$10分 & 论文严密、逻辑性强、表达清楚。 & \\
    \cline{3-5}
    & & 3$\sim$7分 & 基本概念清晰、层次分明。表达较清楚。 & \\
    \cline{3-5}
    & & 3分以下 & 表达较差 & \\
    \hline
    \multicolumn{4}{|r|}{总分} & \\
    \hline
  \end{tabular}

  \vspace{14pt}

  {\zihao{4}\heiti 开题评审小组组成：}

  \vspace{6pt}

  \zihao{5}
  \begin{tabular}{|p{2.5cm}|p{3.5cm}|p{3.0cm}|p{3.6cm}|p{2.6cm}|}
    \hline
    组成 & 姓名 & 职称 & 单位 & 签字 \\
    \hline
    导~师 & & & & \\
    \hline
    组~长 & & & & \\
    \hline
    成~员 & & & & \\
    \hline
    & & & & \\
    \hline
    & & & & \\
    \hline
    & & & & \\
    \hline
  \end{tabular}

  \vspace{10pt}

  年\hspace{2em}月\hspace{2em}日
\end{center}

\newpage

% ==================== 评审意见 ====================
\zihao{4}\heiti 六、评审意见

\vspace{10pt}

\zihao{5}
\begin{tabular}{|p{15.4cm}|}
  \hline
  \textbf{导师对本课题的评价} \\
  \vspace{3.2cm} \\
  导师签名\hspace{4em}\hspace{5.5cm}年\hspace{1.5em}月\hspace{1.5em}日 \\
  \hline
  \textbf{评审小组的审查结论} \\
  \vspace{3.2cm} \\
  组长\hspace{2em}\hspace{6cm}组员\hspace{5cm}年\hspace{1.5em}月\hspace{1.5em}日 \\
  \hline
  \textbf{学科专业分委员会意见} \\
  \vspace{2.4cm} \\
  负责人签名\hspace{3em}\hspace{5.5cm}年\hspace{1.5em}月\hspace{1.5em}日 \\
  \hline
\end{tabular}

\newpage

\end{document}
